\documentclass[11pt]{article}

\usepackage{amsmath,amssymb}
\usepackage{booktabs}
\usepackage{graphicx}
\usepackage{hyperref}
\usepackage[margin=1in]{geometry}
\usepackage{natbib}

\title{Geometric Perturbation Metrics Measure Magnitude, Not Mechanism}

\author{Elliot Tower\\
\texttt{elliot@elliottower.ai}}

\date{}

\begin{document}
\maketitle

\begin{abstract}
Geometric metrics for perturbation biology---direction instability,
Grassmannian geodesic distance, profile variability---are used to
quantify whether a drug or gene knockout acts consistently across
cellular contexts. We demonstrate that variance-based and subspace-based
formulations of these metrics are dominated by effect magnitude. Across
three perturbation modalities (LINCS L1000, Perturb-seq, DepMap CRISPR
screens), we show two mechanistically distinct artifacts: (1) genes with
stronger dependency scores have mechanically higher variance across cell
lines (DepMap: raw Cohen's $d = 2.65$, coefficient-of-variation
normalized $d = -1.33$), and (2) genes with larger transcriptomic
effects produce better-estimated PCA subspaces that shift geodesic
distances regardless of directional diversity (Perturb-seq: raw $d =
0.60$, magnitude-matched $d = -0.06$, $p = 0.62$). A pre-registered
directional prediction that essential gene knockdowns would show greater
divergence than non-essentials technically passed its criterion ($d >
0.35$) on the naive metric; we override this confirmation on honesty
grounds after discovering the magnitude confound. We provide a causal
DAG explaining why the artifact arises, demonstrate that
magnitude-normalized metrics eliminate it, and show that cosine-based
formulations (direction instability) are inherently immune. The
corrective rule: normalize for effect magnitude before comparing groups
defined by effect strength, or geometric ``divergence'' reduces to the
tautology that bigger effects are noisier.
\end{abstract}

\section{Introduction}

High-throughput perturbation screens produce multi-dimensional signatures
for thousands of drugs and gene knockdowns across panels of cell lines.
A natural question follows: does this perturbation act consistently
across contexts, or does it produce divergent effects depending on
cellular background? Geometric metrics formalize this question.
Direction instability measures the mean pairwise angular deviation of
perturbation signatures across cell lines. Grassmannian geodesic
distance measures how much a perturbation's response subspace rotates
between cell types. Profile variability (standard deviation across
contexts) measures raw spread.

These metrics are used to characterize drug mechanism conservation, gene
essentiality, and perturbation quality in LINCS L1000
\citep{subramanian2017next}, Perturb-seq \citep{replogle2022mapping},
JUMP Cell Painting \citep{chandrasekaran2023jump}, and DepMap CRISPR
screens. The implicit assumption is that ``more divergent'' reflects
biology---context-dependent mechanisms, compensatory responses, or
pathway rewiring.

We show this assumption fails for two of the three major metric
formulations. Variance-based metrics (standard deviation of scores
across cell lines) and subspace-based metrics (geodesic distance between
PCA-estimated perturbation subspaces) are confounded by effect
magnitude. Groups with stronger perturbation effects appear more
``divergent'' by construction, regardless of directional diversity.
Cosine-based metrics (direction instability) are immune because they
normalize by magnitude before comparison.

Our evidence comes from a sequence of pre-registered and exploratory
experiments across three data modalities, reported with an explicit
integrity boundary between confirmatory and post-hoc analyses. The
discovery unifies two previously puzzling observations: (1) biological
confound corrections fail to rerank direction instability across three
domains (Spearman $\rho = 0.83$--$0.99$), and (2) essential gene
knockdowns appear to ``diverge'' more than non-essentials ($d = 0.55$ in
Perturb-seq). The first reflects magnitude-dominance of the ranking; the
second reflects magnitude-confounding of the comparison. Both dissolve
under the same DAG.

\section{Metrics and Their Magnitude Sensitivity}

\subsection{Direction instability (cosine-based, magnitude-invariant)}

For a perturbation tested in $K$ contexts with signature vectors
$\mathbf{s}_1, \ldots, \mathbf{s}_K \in \mathbb{R}^d$:
\begin{equation}
  D = 1 - \frac{2}{K(K-1)} \sum_{i < j}
    \frac{\mathbf{s}_i^\top \mathbf{s}_j}
         {\|\mathbf{s}_i\| \, \|\mathbf{s}_j\|}
\end{equation}
The $\ell_2$ normalization makes $D$ invariant to the magnitude of each
$\mathbf{s}_i$. A drug that produces a strong, consistent signal across
cell lines receives the same $D$ as one producing a weak, consistent
signal.

\subsection{Profile variability (variance-based, magnitude-sensitive)}

For a gene with dependency scores $x_1, \ldots, x_K$ across $K$ cell
lines:
\begin{equation}
  V = \text{SD}(x_1, \ldots, x_K)
\end{equation}
This metric grows with the absolute magnitude of effects. A gene with
scores uniformly near $-1$ (strongly essential everywhere) and one with
scores clustered near $0$ (non-essential everywhere) will have different
$V$ even if both are equally ``consistent''---the essential gene's
scores simply occupy a larger range.

\subsection{Grassmannian geodesic distance (subspace-based, indirectly magnitude-sensitive)}

For a gene knocked down in two cell types, let $U_1, U_2 \in
\text{Gr}(k, d)$ denote the top-$k$ PCA subspaces of the perturbation
response. The geodesic distance $\delta(U_1, U_2) =
\|\boldsymbol{\theta}\|_2$ (principal angles) is geometrically
magnitude-invariant. However, $U_1$ and $U_2$ are \emph{estimated} from
finite samples. Perturbations with larger effect sizes produce
higher-SNR data, yielding better-estimated subspaces. This estimation
quality difference propagates into the distance metric, creating an
indirect magnitude sensitivity that is absent from the population-level
definition but present in every finite-sample application.

\section{Pre-registered Experiments}

\subsection{Integrity chain}

All pre-registrations were committed before the author reviewed
experimental output (commit \texttt{249abaf}). The DepMap replication
hypothesis was committed separately (commit \texttt{98897ed}) before
execution. The analysis script was frozen (commit \texttt{92e3276})
before running. The magnitude-confound analysis (Section~\ref{sec:artifact})
is exploratory and labeled as such.

\subsection{Experiment 1: Do biological corrections rerank direction instability?}

\textbf{Hypothesis:} Removing confounding features (toxicity genes in
LINCS, essential-gene subspace in Perturb-seq, cell-health features in
JUMP-CP) would substantially reorganize the direction instability ranking
($\rho < 0.70$).

\textbf{Decision criteria:} $\rho < 0.70$: correction bites.
$0.70 \leq \rho < 0.83$: marginal. $\rho \geq 0.83$: does not bite.

\textbf{Result:} Corrections do not bite in any domain (Table~\ref{tab:robustness}).

\begin{table}[h]
\centering
\caption{Rank preservation after biological corrections across three
perturbation domains. All Spearman correlations exceed the pre-registered
``doesn't bite'' threshold of $\rho \geq 0.83$.}
\label{tab:robustness}
\begin{tabular}{llcc}
\toprule
Domain & Correction & $\rho$ & $N$ \\
\midrule
LINCS (global) & Toxicity genes removed & $>0.99$ & 8{,}949 \\
LINCS (HDAC gradient) & Toxicity genes removed & $0.83$ & 20 \\
Perturb-seq & Essential-gene subspace removed & $0.91$ & 1{,}676 \\
JUMP-CP & Cell-health features removed & $0.99$ & 25{,}254 \\
\midrule
Same-tissue LINCS & Within-tissue variance genes & $0.98$--$0.99$ & 176--358 \\
\bottomrule
\end{tabular}
\end{table}

\subsection{Experiment 2: Does essential-gene divergence replicate in DepMap?}

\textbf{Exploratory prior (disclosed):} In Perturb-seq, essential gene
knockdowns showed higher Grassmannian geodesic distance than non-essentials
(Cohen's $d = 0.55$, bootstrap CI $[0.45, 0.65]$, $p < 10^{-30}$).

\textbf{Pre-registered prediction:} Essential genes will show higher
variability of dependency profiles across cell lines than non-essential
genes (one-tailed, $d > 0.35$). Framed as confirmatory replication of
the Perturb-seq finding.

\textbf{Operationalization:} Standard deviation of PC1-corrected Chronos
scores across 1{,}178 cell lines. Essential genes defined by DepMap
common-essentials list ($n = 1{,}242$); non-essential genes defined by
DepMap nonessentials list ($n = 726$).

\textbf{Result:} $d = 2.65$ (bootstrap CI $[2.52, 2.79]$, $p <
10^{-297}$). Pre-registered criterion satisfied. PC1 removal changes
effect by 0.2\%, ruling out the fitness confound.

\section{Discovery of the Magnitude Artifact}
\label{sec:artifact}

\textbf{Status: EXPLORATORY. Not pre-registered. Discovered post-hoc
after observing that $d = 2.65$ was implausibly large.}

The $d = 2.65$ result triggered scrutiny: Cohen's $d$ values above 2
imply near-zero distributional overlap, which is unusual for biological
comparisons. Two checks revealed the effect is a magnitude artifact.

\subsection{DepMap: coefficient-of-variation normalization}

Essential genes have large negative Chronos scores (by definition---this
is what ``essential'' means in DepMap). Larger absolute scores
mechanically produce larger absolute variance. Normalizing by mean
effect magnitude (coefficient of variation: SD / $|\bar{x}|$) removes
this confound:

\begin{itemize}
  \item Essential CV: mean $= 0.38$, median $= 0.26$
  \item Non-essential CV: mean $= 2.95$, median $= 1.71$
  \item Cohen's $d$ (CV): $-1.33$ (sign reversal)
\end{itemize}

After normalization, non-essential genes show \emph{higher} relative
variability. The raw $d = 2.65$ was the tautology that ``things with
bigger effects have bigger variance.''

\subsection{Perturb-seq: magnitude-matched comparison}

Geodesic distance correlates with CERES effect magnitude (Spearman $\rho
= 0.33$, $p < 10^{-40}$). After matching essential and non-essential
genes by response magnitude ($n = 127$ matched pairs within 0.2 CERES
units):

\begin{itemize}
  \item Essential mean distance: 2.659
  \item Non-essential mean distance: 2.671
  \item Cohen's $d$: $-0.06$ ($p = 0.62$, Mann-Whitney)
\end{itemize}

The effect vanishes entirely. The original $d = 0.55$ was driven by
essential genes having larger transcriptomic effects, which improves PCA
subspace estimation quality and shifts geodesic distances independent of
true directional diversity.

\subsection{Overriding the pre-registered confirmation}

The pre-registered criterion ($d > 0.35$) was technically satisfied by
the raw SD metric ($d = 2.65$). We override this on honesty grounds: the
metric captured magnitude, not the directional divergence it was
intended to measure. A pre-registration protects against cherry-picking
hypotheses; it does not oblige one to report a confounded metric as if
it measures the intended construct. Both the ``confirmation'' and the
artifact are reported; the latter is our actual conclusion.

\section{Causal DAG}
\label{sec:dag}

The magnitude confound has the same causal structure across both
affected metrics, though the mechanism differs.

\subsection{DepMap (variance pathway)}

\begin{equation*}
\text{Magnitude} \rightarrow
\begin{cases}
\text{Essentiality label} & \text{(confounds grouping)} \\
\text{Score variance} & \text{(inflates metric)}
\end{cases}
\end{equation*}

Essentiality is \emph{operationally defined} by dependency-score
magnitude (Chronos $< -0.5$). The magnitude--essentiality association is
near-tautological. Magnitude also causes higher variance (larger
absolute numbers span larger ranges). Adjusting for magnitude is correct
because it removes a spurious path, not a mediating one.

\subsection{Perturb-seq (estimation pathway)}

\begin{equation*}
\text{Magnitude} \rightarrow
\begin{cases}
\text{Essentiality status} & \text{(confounds grouping)} \\
\text{Subspace estimation quality} \rightarrow \text{Geodesic distance}
  & \text{(inflates metric)}
\end{cases}
\end{equation*}

Here magnitude is not tautologically linked to essentiality, but the
confounding path through estimation quality produces the same artifact.
The evidence that this is a confounder rather than a mediator: if
magnitude were a \emph{mediator} (essential $\rightarrow$ big response
$\rightarrow$ genuinely different direction), then conditioning on
magnitude would attenuate but not eliminate the essential/non-essential
difference. Instead, magnitude-matching eliminates it completely ($d =
-0.06$). The entire signal traversed the estimation-quality path.

\subsection{LINCS direction instability (no confound)}

\begin{equation*}
\text{Magnitude} \not\rightarrow D \quad \text{(cosine normalization
blocks the path)}
\end{equation*}

Direction instability divides by $\|\mathbf{s}_i\|$ before computing
angles. Magnitude cannot propagate to the metric. This explains the
observed robustness: corrections do not rerank because the metric
measures direction independent of magnitude, and the direction is
stable.

\section{The Fix}

Three normalizations remove the magnitude confound from affected
metrics:

\begin{enumerate}
  \item \textbf{Coefficient of variation} (SD / $|\text{mean effect}|$):
    removes scale dependence. Applicable when effects have non-zero means.
  \item \textbf{Magnitude matching}: compare only perturbations with
    similar absolute effect sizes. The most assumption-free approach; does
    not require linearity.
  \item \textbf{Residualization}: regress out effect magnitude and test
    residuals. Assumes a linear (or specified) relationship between
    magnitude and the metric.
\end{enumerate}

Cosine-based formulations (direction instability, Equation~1) require no
correction---they are the magnitude-normalized version by construction.

\section{Discussion}

This paper makes one claim: variance-based and subspace-based
perturbation-divergence metrics are magnitude-confounded when used to
compare groups defined by effect strength. The claim is supported by two
mechanistically distinct artifacts (variance inflation in DepMap,
estimation-quality bias in Perturb-seq) that converge on the same
conclusion across three data modalities.

The practical implication is a rule: before comparing ``divergence''
across groups that differ in perturbation strength, normalize for
magnitude. Failing to do so produces effects that look large ($d = 2.65$,
$d = 0.55$) but reduce to the tautology that bigger effects are noisier
or better-estimated.

\subsection{Limitations}

The magnitude-matching analysis in Perturb-seq used $n = 127$ matched
pairs (constrained by the overlap in effect-size distributions between
essential and non-essential genes). Larger matched samples from future
screens would tighten confidence intervals around the null.

Our analysis does not address whether magnitude normalization could mask
\emph{real} directional divergence in domains where essentiality is not
defined by magnitude. In principle, a biological process could produce
both larger effects AND genuinely different directions. The DAG
distinguishes these cases (mediator vs.\ confounder), but the
distinction requires per-domain justification.

\subsection{What the robustness result means in light of the artifact}

The bracket-norm robustness finding (corrections do not rerank, $\rho =
0.83$--$0.99$) is independently valid because direction instability is
cosine-based and magnitude-invariant. The robustness reflects genuine
directional stability of perturbation signatures, not magnitude
dominance. However, the \emph{interpretation} shifts: the ranking is
robust because the underlying geometry is stable, which we now
understand is partially because cosine normalization removes the primary
source of apparent variation (magnitude).

\subsection{Implications for the field}

Researchers computing variability metrics on perturbation atlases
(DepMap, Perturb-seq, PRISM, JUMP-CP) should verify that their metric is
magnitude-invariant before interpreting group differences as biological
divergence. Metrics that lack built-in normalization (SD, geodesic
distance on estimated subspaces) require explicit magnitude correction.
Metrics with built-in normalization (cosine similarity, direction
instability) are safe by construction.

\section{Conclusion}

We pre-registered a divergence prediction, the naive metric confirmed
it, and our own post-hoc checks killed it. The killing is the finding:
geometric perturbation metrics measure magnitude unless explicitly
prevented from doing so. The DAG is simple, the fix is simple, and the
failure mode is likely widespread in any analysis comparing perturbation
groups that differ in effect strength.

\bibliographystyle{plainnat}
\bibliography{references}

\end{document}
